% page 192 du fac-simile (image p-191 du scan)
% bloc 195.6 x 254.9 mm ; marge gauche 3.0 mm
\pgc{-17.780mm}{0.000mm}{
\l{}
\l{\cel{18}184}
\l{}
\l{}
\l{}
\l{\cel{13}\sou{furn}azo. Forno inkandecanta. (\sou{anke metaf.)}}
\l{}
\l{\cel{13}\sou{furnisar}. (\sou{trans}.)\cel{2}Livrar sufico-quante, sive vende, sive}
\l{\cel{16}donace. - EFI.}
\l{}
\l{\cel{13}\sou{furtar}.\cel{2}(\sou{trans.}) Proprigar a su, per ruzo o per violento, to}
\l{\cel{16}quo savate esas la proprajo, la proprietajo, di altru.-ISL.}
\l{}
\l{\cel{13}\sou{furunklo}. (\sou{patol.}) Tumoro inflamala limitoza, kun saliajo en}
\l{\cel{16}lua centro, qua tempo-finas per pusifo. - DEFIS.}
\l{}
\l{\cel{13}\sou{fushar}. (\sou{trans.)}\cel{2}Exekutar tasko neglijoze. - D.}
\l{}
\l{\cel{13}\sou{fusilo}. Paf-armo portebla, ek fer-tubo muntita sur ligna stango}
\l{\cel{16}krosoza, provizita ye mekanismo di qua la deskuplo acendas}
\l{\cel{16}la kargajo pulvero di la armo. - FIRS.}
\l{}
\l{\cel{13}\sou{fusto}. I. (\sou{arkitekt.}) Stango di kolono, inter la bazo e la ka-}
\l{\cel{16}pitelo. - II. (\sou{artilrio}). Suportilo, di kanono, di lorneto,}
\l{\cel{16}di teleskopo, e c. - EFIS.}
\l{}
\l{\cel{13}\sou{fusteno}. Stofo di qua la fili esas krucumita - la fili longes-}
\l{\cel{16}ala ek lino, la fili transversa ek kotono, uzata por facar}
\l{\cel{16}subjupi, kamizoli, e c. - EFIS.}
\l{}
\l{\cel{13}\sou{futero}. Stofo ye qua onu garnisas la parto interna di vesti.}
\l{\cel{16}- DIRS.}
\l{}
\l{\cel{13}\sou{futuro}. I. (\sou{gram.)}\cel{2}Tempo di verbo, qua expresas to quo esos.}
\l{\cel{16}- II. To quo esas eventonta, existonta. - DEF.}
\l{}
\l{\cel{13}\sou{fuxio}. (\sou{bot.}) Orno-planto di qua la flori esas pendanta, genero}
\l{\cel{16}de la familio \textquotedbl{}enoterei\textquotedbl{}. - L.\cel{1}\sou{fuchsia}.}
\l{}
\l{\cel{13}\sou{fuxi}no. (\sou{kemio)}\cel{2}Farbo reda, preparita ek anilino. -\cel{1}\sou{Simbolo}}
\l{\cel{16}\sou{kemiala}\cel{1}: C20\cel{2}H20\cel{2}N3Cl,\cel{2}E\cel{2}H2\cel{2}O.}
\l{}
\l{\cel{13}\sou{fuzar}. (\sou{trans}.)Liquidigar (korpo solida) per la efiko di kaloro.}
\l{\cel{16}- EFIS.}
\l{}
\l{\cel{13}\sou{fuzeo}. Pirotekno-kozo inkluzita en kartono (kartocho) qua}
\l{\cel{16}acendesas stratope, e lansas spricaji ek partikuli kombus-}
\l{\cel{16}tata. - F.}
\l{}
\l{\cel{13}\sou{fuzelo}. (\sou{geom.}) Parto di surfaco sferala, inter du duima granda}
\l{\cel{16}cirkli - di surfaco cilindrala o konala, inter du plani}
\l{\cel{16}facita segun la axo. - F.}
\l{}
\l{\cel{13}\sou{fuzeno}. (\sou{bot}.) Arboreto di qua la ligno uzesas por fabrikar}
\l{\cel{16}spindeli lardo-spici, e c., e di qua la rameti karbonigita}
\l{\cel{16}donas krayoni lejera por skisado. - L.\cel{1}\sou{evonymus}. - FIS.}
\l{}
\l{\cel{13}\sou{fyordo}. Singla de la sesguri streta e profunda qui festonizas}
\l{\cel{16}la litoro westala di Skandinavia e di Skotia. - DEFIRS.}
}
